%% Lecture 06 — Covariant and Contravariant Tensors; Relativistic Mechanics
\section{Lecture 06 — Covariant and Contravariant Tensors; Relativistic Mechanics}\label{sec:lec06}

\begin{center}
\begin{tabular}{ll}
  \textbf{Date:}\quad 21/04/2026   & \\
\end{tabular}
\end{center}
\hrule\vspace{1.5em}

\subsection*{Overview}
This lecture has two parts.  In the first part we complete the tensor
formalism begun in Lecture~05.  We introduce the distinction between
\textbf{contravariant} (upper-index) and \textbf{covariant} (lower-index)
objects, derive the transformation law for covariant four-vectors from first
principles, discuss mixed-rank tensors and why the position of each index
matters, and prove that the Minkowski metric is itself a Lorentz-invariant
rank-2 tensor.  We also give a self-contained proof that any full contraction
of tensor indices yields a Lorentz scalar.

In the second part we build \textbf{relativistic mechanics} from the action
principle.  Following the modern approach—identify the symmetries first,
then write the most general action consistent with them—we show that the
symmetries of special relativity essentially fix the free-particle action
uniquely.  Matching the non-relativistic limit fixes the remaining constant,
and the Euler--Lagrange equations then deliver the relativistic momentum,
energy, and dispersion relation.  We close by adding the minimal
electromagnetic coupling and recovering the covariant Lorentz force law.

% ==============================================================================
\subsection{Tensors in Special Relativity (Concluded)}
% ==============================================================================

\subsubsection*{Recap: Contravariant Four-Vectors}

In the previous lecture we defined a \textbf{four-vector} as any set of four
components $V^\mu$ ($\mu = 0,1,2,3$) that transforms under a Lorentz
transformation in exactly the same way as the coordinate differential
$dx^\mu$:
\begin{equation}\label{eq:lec06:contravariant-transform}
    V'^\mu = \Lambda^\mu_{\ \nu}\,V^\nu .
\end{equation}
An object satisfying this law is called a \textbf{contravariant four-vector};
its single index is written \emph{upstairs}.  Not every collection of four
numbers is a four-vector: the transformation law must be verified.

Standard examples: the position four-vector $x^\mu = (ct,\bvec{x})$, the
four-velocity $u^\mu = \gamma(c,\bvec{v})$, the four-momentum $p^\mu$, and
the four-potential $A^\mu$.

\subsubsection*{Covariant Four-Vectors: Lowering an Index}

Given a contravariant four-vector $V^\mu$, we define a companion object with
a \emph{lower} index by contracting with the Minkowski metric
$\eta_{\mu\nu} = \mathrm{diag}(+1,-1,-1,-1)$:
\begin{equation}\label{eq:lec06:covariant-def}
    V_\mu \;\equiv\; \eta_{\mu\nu}\,V^\nu .
\end{equation}
Because $\eta_{\mu\nu}$ is diagonal with entries $(+1,-1,-1,-1)$, the
explicit components are
\begin{equation}\label{eq:lec06:covariant-components}
    V_\mu = \bigl(V^0,\;-V^1,\;-V^2,\;-V^3\bigr).
\end{equation}
This object $V_\mu$ is called a \textbf{covariant four-vector} (or
\textbf{covector}).  It is a genuinely new object: $V_\mu \neq V^\mu$ because
the Minkowski metric is not the identity.  We keep the same letter but use a
\emph{downstairs} index to distinguish the two.

\subsubsection*{Transformation Law of the Covariant Four-Vector}

We derive how $V_\mu$ transforms.  Starting from the definition in the primed
frame:
\begin{align}
    V'_\mu
    &= \eta_{\mu\nu}\,V'^\nu
     = \eta_{\mu\nu}\,\Lambda^\nu_{\ \rho}\,V^\rho .
     \label{eq:lec06:cov-step1}
\end{align}
We want the right-hand side expressed in terms of $V_\sigma$ (lower index),
not $V^\rho$ (upper index).  Insert $V^\rho = \eta^{\rho\sigma}V_\sigma$
(raising with the inverse metric $\eta^{\mu\nu}$, satisfying
$\eta_{\mu\rho}\eta^{\rho\nu}=\delta^\nu_\mu$):
\begin{equation}\label{eq:lec06:cov-step2}
    V'_\mu
    = \underbrace{\eta_{\mu\nu}\,\Lambda^\nu_{\ \rho}\,\eta^{\rho\sigma}}_{\displaystyle\equiv\;\tilde\Lambda_\mu^{\ \sigma}}\;
      V_\sigma .
\end{equation}
We define the \textbf{covariant Lorentz matrix}
\begin{equation}\label{eq:lec06:lambda-covariant}
    \tilde\Lambda_\mu^{\ \sigma}
    \;\equiv\;
    \eta_{\mu\nu}\,\Lambda^\nu_{\ \rho}\,\eta^{\rho\sigma} .
\end{equation}

\textbf{Claim:} $\tilde\Lambda_\mu^{\ \sigma} = (\Lambda^{-1})^\sigma_{\ \mu}$.

\textit{Proof.}  Lorentz transformations are defined by the condition that
they preserve the Minkowski metric:
\begin{equation}\label{eq:lec06:lorentz-condition}
    \eta_{\alpha\beta}\,\Lambda^\alpha_{\ \mu}\,\Lambda^\beta_{\ \nu}
    = \eta_{\mu\nu}
    \qquad\Longleftrightarrow\qquad
    \Lambda^T\eta\,\Lambda = \eta .
\end{equation}
Rearranging: $\Lambda^{-1} = \eta^{-1}\Lambda^T\eta = \eta\Lambda^T\eta$
(since $\eta^{-1}=\eta$ numerically in flat space).  In components:
\begin{equation}\label{eq:lec06:lambda-inverse-proof}
    (\Lambda^{-1})^\sigma_{\ \mu}
    = \eta^{\sigma\rho}\,\Lambda^\nu_{\ \rho}\,\eta_{\nu\mu}
    = \eta_{\mu\nu}\,\Lambda^\nu_{\ \rho}\,\eta^{\rho\sigma}
    = \tilde\Lambda_\mu^{\ \sigma} .
    \qquad \square
\end{equation}
Therefore the transformation law of a covariant four-vector is
\begin{equation}\label{eq:lec06:covariant-final}
    \boxed{V'_\mu = (\Lambda^{-1})^\nu_{\ \mu}\,V_\nu .}
\end{equation}
A contravariant four-vector transforms with $\Lambda$; a covariant four-vector
transforms with the inverse $\Lambda^{-1}$.  Both are legitimate
``four-vectors'' — they are distinguished only by the position of the index
and by their respective transformation matrices.

\subsubsection*{The Inner Product is a Lorentz Scalar}

Contracting one contravariant and one covariant four-vector:
\begin{equation}\label{eq:lec06:inner-product}
    V^\mu W_\mu
    = V^0 W^0 - V^1 W^1 - V^2 W^2 - V^3 W^3 ,
\end{equation}
is automatically a Lorentz scalar.  Under a Lorentz transformation:
\begin{equation}\label{eq:lec06:inner-product-proof}
    V'^\mu W'_\mu
    = \bigl(\Lambda^\mu_{\ \alpha}V^\alpha\bigr)
      \bigl((\Lambda^{-1})^\beta_{\ \mu}W_\beta\bigr)
    = (\Lambda^{-1})^\beta_{\ \mu}\,\Lambda^\mu_{\ \alpha}\;V^\alpha W_\beta
    = \delta^\beta_{\ \alpha}\,V^\alpha W_\beta
    = V^\alpha W_\alpha .
\end{equation}
The product of the two Lorentz matrices is the identity by definition.  This
illustrates the core bookkeeping rule: \emph{any fully contracted expression
— with every free index appearing once upstairs and once downstairs — is
Lorentz invariant.}

\subsubsection*{Classification of Tensor Objects}

\begin{itemize}
    \item \textbf{Scalar (rank 0):} no free indices; invariant under all
      Lorentz transformations.  Examples: $ds^2$, $m$, $d\tau$,
      $V^\mu W_\mu$.
    \item \textbf{Contravariant four-vector (rank 1, upper):}
      $V'^\mu = \Lambda^\mu_{\ \nu}V^\nu$.
    \item \textbf{Covariant four-vector (rank 1, lower):}
      $V'_\mu = (\Lambda^{-1})^\nu_{\ \mu}V_\nu$.
    \item \textbf{Rank-2 contravariant tensor:}
      $T'^{\mu\nu} = \Lambda^\mu_{\ \rho}\Lambda^\nu_{\ \sigma}T^{\rho\sigma}$.
    \item \textbf{Rank-2 covariant tensor:}
      $T'_{\mu\nu}
      = (\Lambda^{-1})^\rho_{\ \mu}(\Lambda^{-1})^\sigma_{\ \nu}T_{\rho\sigma}$.
    \item \textbf{Mixed rank-2 tensor:}
      $T'^\mu_{\ \nu} = \Lambda^\mu_{\ \rho}(\Lambda^{-1})^\sigma_{\ \nu}T^\rho_{\ \sigma}$.
    \item \textbf{Rank-$n$ tensor:} each upper index transforms with a factor
      $\Lambda$; each lower index transforms with a factor $\Lambda^{-1}$.
\end{itemize}

\subsubsection*{Index Position Matters in Mixed Tensors}

For a tensor with at least one upper and one lower index, the \emph{relative
position} of the indices matters.  In general:
\begin{equation}\label{eq:lec06:index-position}
    T^\mu_{\ \nu} \;\neq\; T_{\ \nu}^\mu ,
\end{equation}
because these represent different index-lowering choices applied to a tensor
that may not be symmetric.  When writing mixed tensors, always specify which
index was raised or lowered and preserve its left-right order relative to the
other indices.  For example, $T^\mu_{\ \nu}$ has the upper index to the left,
while $T_{\ \nu}^\mu$ has the upper index to the right; unless $T$ is
symmetric these are different components.

\subsubsection*{Raising and Lowering Indices}

For any tensor, any index can be raised or lowered independently using the
metric or its inverse:
\begin{align}
    V_\mu &= \eta_{\mu\nu}\,V^\nu, &
    V^\mu &= \eta^{\mu\nu}\,V_\nu, \label{eq:lec06:raise-lower-vector} \\[4pt]
    T_{\mu\nu} &= \eta_{\mu\rho}\,T^\rho_{\ \nu}
               = \eta_{\mu\rho}\,\eta_{\nu\sigma}\,T^{\rho\sigma},
    \label{eq:lec06:lower-rank2}
\end{align}
and analogously for higher-rank tensors.  Applying $\eta$ followed by
$\eta^{-1}$ (or vice versa) returns the original tensor.

With the $(+,-,-,-)$ metric, lowering a time-component index leaves it
unchanged while lowering a spatial-component index introduces a minus sign:
\begin{equation}\label{eq:lec06:lowering-sign}
    V_0 = V^0, \qquad V_i = -V^i \quad (i = 1,2,3).
\end{equation}

\subsubsection*{Contraction and the Trace}

\textbf{Contraction} means setting one upper and one lower index equal and
summing (Einstein convention).  Each contraction reduces the tensor rank
by~2:

\begin{itemize}
    \item A rank-2 tensor $T^{\mu\nu}$ contracted to a scalar:
\begin{equation}\label{eq:lec06:trace-rank2}
    T \equiv T^\mu_{\ \mu} = \eta_{\mu\nu}\,T^{\mu\nu} .
\end{equation}
    \item A rank-3 tensor $T^{\mu\nu\rho}$ contracted on the first and third
      indices to a four-vector:
\begin{equation}\label{eq:lec06:contract-rank3}
    S^\nu = T^\mu_{\ \nu\mu} = \eta_{\mu\rho}\,T^{\mu\nu\rho} .
\end{equation}
    \item More generally, contracting $k$ pairs of indices in a rank-$n$
      tensor yields a rank-$(n-2k)$ tensor.
\end{itemize}

\subsubsection*{Why Every Full Contraction is a Lorentz Scalar}

We verify directly for the trace of a rank-2 tensor.  In the primed frame:
\begin{align}
    T'
    &= \eta_{\mu\nu}\,T'^{\mu\nu} \notag \\
    &= \eta_{\mu\nu}\,\Lambda^\mu_{\ \rho}\,\Lambda^\nu_{\ \sigma}\,T^{\rho\sigma} \notag \\
    &= T^{\rho\sigma}\bigl(\eta_{\mu\nu}\,\Lambda^\mu_{\ \rho}\,\Lambda^\nu_{\ \sigma}\bigr) \notag \\
    &= T^{\rho\sigma}\,\eta_{\rho\sigma} = T .
    \label{eq:lec06:scalar-proof}
\end{align}
In the third step we rearranged numerical factors (permitted since sums of
products of numbers commute).  In the fourth step we used the defining
property~\eqref{eq:lec06:lorentz-condition} of Lorentz transformations.
The same argument extends to any rank: each pair of contracted indices
produces a factor $(\Lambda^{-1})\Lambda = \mathbf{1}$.

\subsubsection*{The Minkowski Metric is a Lorentz-Invariant Tensor}

We now show that $\eta_{\mu\nu}$ is itself a rank-2 covariant tensor — and a
\emph{constant} one: it is numerically the same in every inertial frame.

Applying the covariant rank-2 transformation law to $\eta_{\mu\nu}$:
\begin{align}
    \eta'_{\mu\nu}
    &= (\Lambda^{-1})^\rho_{\ \mu}\,(\Lambda^{-1})^\sigma_{\ \nu}\,\eta_{\rho\sigma} .
    \label{eq:lec06:metric-transform}
\end{align}
Inserting $(\Lambda^{-1})^\rho_{\ \mu} = \eta_{\mu\alpha}\Lambda^\alpha_{\ \beta}\eta^{\beta\rho}$
and collecting, the right-hand side reduces to $\eta_{\mu\nu}$ by the
Lorentz condition~\eqref{eq:lec06:lorentz-condition}:
\begin{equation}\label{eq:lec06:metric-invariant}
    \boxed{\eta'_{\mu\nu} = \eta_{\mu\nu} .}
\end{equation}
Equation~\eqref{eq:lec06:metric-invariant} is not a coincidence: it is exactly the
defining property of Lorentz transformations.  The Minkowski metric is therefore a
\textbf{constant tensor}: it looks the same to every inertial observer.

This completes the tensor vocabulary required for the rest of the course.

% ==============================================================================
\subsection{Relativistic Mechanics}
% ==============================================================================

\subsubsection*{The Modern Paradigm: Symmetries Determine the Action}

The Lagrangian/action formulation of mechanics says that a physical system is
specified by its action $S = \int L\,dt$.  Rather than guessing $L$, the
modern approach proceeds in three steps:
\begin{enumerate}
    \item \textbf{Identify the symmetries.}  For a relativistic point particle
      in flat spacetime these are: Lorentz boosts, spatial rotations,
      spatial-translation invariance (homogeneity of space), and
      time-translation invariance (homogeneity of time).
    \item \textbf{Write the most general action consistent with those
      symmetries.}  Any term that is not a Lorentz scalar violates Lorentz
      symmetry and is forbidden.  Spatial/temporal homogeneity forbids
      explicit dependence on $x^\mu$ itself.
    \item \textbf{Fix any remaining undetermined constants} by demanding
      agreement with the known non-relativistic limit.
\end{enumerate}
This symmetry-first approach is how all of modern field theory is built.  For
a free particle the constraints are so tight that the action is essentially
unique.

\subsubsection*{The Only Available Lorentz Scalar for a Free Particle}

For a free point particle the dynamical degrees of freedom are its worldline
coordinates $x^\mu(\lambda)$ and the only dimensional constant is $c$.  We
must construct a Lorentz scalar that:
\begin{itemize}
    \item depends only on $dx^\mu$ (translational invariance forbids $x^\mu$
      itself),
    \item is reparameterisation-invariant (independent of the labelling
      parameter $\lambda$).
\end{itemize}
The unique (up to a constant factor) object satisfying all requirements is the
spacetime line element:
\begin{equation}\label{eq:lec06:line-element}
    ds^2 \;\equiv\; \eta_{\mu\nu}\,dx^\mu\,dx^\nu
         = c^2\,dt^2 - d\vec{x}^{\,2}
         > 0
    \qquad(\text{timelike worldline}),
\end{equation}
giving the invariant arc length $ds = c\,dt\sqrt{1-v^2/c^2} = c\,d\tau$.
Any polynomial in $x^\mu$ would violate translational invariance; higher-order
differential invariants reduce to powers of $ds^2$ and can be absorbed into
the constant.  The action must therefore take the form
\begin{equation}\label{eq:lec06:action-candidate}
    S = -\alpha \int ds
\end{equation}
for some constant $\alpha$ to be determined.

\subsubsection*{Why the Minus Sign?}

In Minkowski space, a straight (inertial) worldline \emph{maximises} the
proper time $\tau = \int d\tau = \int ds/c$ between two fixed events.  (This
is the geometrical content of the twin paradox: any deviation from a straight
worldline shortens the elapsed proper time.)  Therefore $\int ds$ is
\emph{maximised} along the classical trajectory, not minimised.

The principle of \textbf{least action} demands $\delta S = 0$ with $S$
stationary at a \emph{minimum}.  Inserting the minus sign converts the
maximum of $\int ds$ into a minimum of $S = -\alpha\int ds$ (for $\alpha > 0$),
restoring compatibility with the variational principle.

\subsubsection*{Fixing the Constant \texorpdfstring{$\alpha$}{$\alpha$}: the Non-Relativistic Limit}

Substituting $ds = c\,dt\sqrt{1-\beta^2}$ (where $\beta = v/c$):
\begin{equation}\label{eq:lec06:action-expanded}
    S = -\alpha c \int \sqrt{1-\beta^2}\,dt .
\end{equation}
Expanding for $v \ll c$:
\begin{equation}\label{eq:lec06:sqrt-expansion}
    \sqrt{1-\beta^2} \approx 1 - \frac{v^2}{2c^2} + O(v^4/c^4),
\end{equation}
so
\begin{equation}\label{eq:lec06:action-nr-expand}
    S \approx -\alpha c \int dt + \frac{\alpha}{2c}\int v^2\,dt .
\end{equation}
The first term is a constant added to $L$ (it does not affect the equations
of motion).  Matching the second term to the classical kinetic action
$\frac{1}{2}m\int v^2\,dt$ requires
\begin{equation}\label{eq:lec06:alpha-fixed}
    \frac{\alpha}{2c} = \frac{m}{2}
    \qquad\Longrightarrow\qquad
    \alpha = mc .
\end{equation}
The free relativistic particle action and its Lagrangian are therefore:
\begin{align}
    S &= -mc\int ds = -mc^2\int\sqrt{1-\frac{v^2}{c^2}}\,dt, \label{eq:lec06:action-final}\\[4pt]
    L &= -mc^2\sqrt{1-\frac{v^2}{c^2}} = -\frac{mc^2}{\gamma}. \label{eq:lec06:lagrangian-final}
\end{align}

\subsubsection*{Non-Relativistic Limit}

Expanding~\eqref{eq:lec06:action-final}:
\begin{equation}\label{eq:lec06:action-nr-check}
    S \approx -mc^2\int dt + \frac{1}{2}m\int v^2\,dt .
\end{equation}
The constant $-mc^2\int dt$ shifts $L$ by a constant and is irrelevant for
the dynamics.  The remaining term $\frac{1}{2}mv^2$ is exactly the classical
free-particle Lagrangian, confirming internal consistency.

\subsubsection*{Conservation Laws from Noether's Theorem}

The Lagrangian~\eqref{eq:lec06:lagrangian-final} depends on $v = |\dot{\vec{x}}|$
but neither on the coordinates $x^i$ nor on $t$ explicitly:
\begin{equation}\label{eq:lec06:noether-symmetries}
    \frac{\partial L}{\partial x^i} = 0, \qquad \frac{\partial L}{\partial t} = 0 .
\end{equation}
By Noether's theorem:
\begin{itemize}
    \item \textbf{Spatial homogeneity} ($\partial L/\partial x^i = 0$)
      $\Rightarrow$ the canonical momentum $p_i = \partial L/\partial \dot x^i$
      is a constant of the motion (conserved for all time).
    \item \textbf{Temporal homogeneity} ($\partial L/\partial t = 0$)
      $\Rightarrow$ the energy $E = \dot x^i\,(\partial L/\partial \dot x^i) - L$
      is conserved.
\end{itemize}
We now compute both.

\subsubsection*{Relativistic Momentum}

\begin{equation}\label{eq:lec06:momentum-calc}
    p_i = \frac{\partial L}{\partial v^i}
        = \frac{\partial}{\partial v^i}\!\left(-mc^2\sqrt{1-\frac{v^2}{c^2}}\right)
        = \frac{-mc^2 \cdot (-v^i/c^2)}{\sqrt{1-v^2/c^2}}
        = \frac{m v^i}{\sqrt{1-v^2/c^2}} = \gamma m v^i .
\end{equation}
In vector form:
\begin{equation}\label{eq:lec06:momentum}
    \boxed{\vec{p} = \gamma m \vec{v}} .
\end{equation}
For $v \ll c$, $\gamma \to 1$ and $\vec{p} \to m\vec{v}$, recovering Newton.
For $v \to c$, $\gamma \to \infty$ and no finite force can accelerate the
particle to the speed of light.

\subsubsection*{Relativistic Energy}

\begin{align}
    E &= \vec{v}\cdot\frac{\partial L}{\partial \vec{v}} - L
       = \vec{v}\cdot\vec{p} - L \notag \\
      &= \gamma m v^2 + mc^2\sqrt{1-\beta^2}
       = \gamma m v^2 + \frac{mc^2}{\gamma} .
       \label{eq:lec06:energy-step}
\end{align}
Using $v^2 = c^2(1 - \gamma^{-2})$:
\begin{equation}\label{eq:lec06:energy-simplify}
    E = \gamma m c^2\!\left(1 - \frac{1}{\gamma^2}\right) + \frac{mc^2}{\gamma}
      = \gamma mc^2 - \frac{mc^2}{\gamma} + \frac{mc^2}{\gamma}
      = \gamma mc^2 .
\end{equation}
Therefore:
\begin{equation}\label{eq:lec06:energy}
    \boxed{E = \gamma mc^2} .
\end{equation}
Expanding for $v \ll c$:
\begin{equation}\label{eq:lec06:energy-nr}
    E = mc^2\!\left(1 + \frac{v^2}{2c^2} + \cdots\right)
      = mc^2 + \frac{1}{2}mv^2 + \cdots
\end{equation}
The energy contains the \textbf{rest energy} $mc^2$ plus the classical kinetic
energy as the leading correction.  The rest energy is unobservable in purely
mechanical processes but becomes physically relevant whenever mass is
converted into radiation or kinetic energy.

\subsubsection*{Energy-Momentum Dispersion Relation}

From~\eqref{eq:lec06:momentum} and~\eqref{eq:lec06:energy} one can eliminate
$\gamma$ and $v$ to obtain a frame-independent relation between $E$ and
$|\vec{p}| \equiv p$:
\begin{align}
    E^2 - p^2c^2
    &= (\gamma mc^2)^2 - (\gamma mv)^2c^2
     = \gamma^2 m^2c^2(c^2 - v^2)
     = \gamma^2 m^2c^4\!\left(1-\frac{v^2}{c^2}\right)
     = m^2c^4 .
     \label{eq:lec06:dispersion-derive}
\end{align}
\begin{equation}\label{eq:lec06:dispersion}
    \boxed{E^2 = p^2c^2 + m^2c^4} .
\end{equation}
This relation is more fundamental than the individual formulas for $E$ and
$\vec{p}$: the combination $E^2 - p^2c^2$ is a Lorentz scalar equal to
$m^2c^4$.  The four-momentum $p^\mu = (E/c,\,\vec{p})$ is a contravariant
four-vector, and~\eqref{eq:lec06:dispersion} is the statement
$\eta_{\mu\nu}p^\mu p^\nu = p^\mu p_\mu = m^2c^2$ (in the $(+,-,-,-)$
convention).

\subsubsection*{Massless Particles Must Travel at the Speed of Light}

Setting $m = 0$ in~\eqref{eq:lec06:dispersion}:
\begin{equation}\label{eq:lec06:massless}
    E = pc .
\end{equation}
From $\vec{p} = \gamma m \vec{v}$ and $E = \gamma mc^2$, the ratio gives
$v = pc^2/E = c$ as long as $m\neq 0$.  Taking $m\to 0$ while keeping
$E = pc$ finite yields $v = c$.  Therefore:
\begin{equation}\label{eq:lec06:massless-speed}
    \boxed{m = 0 \;\;\Longrightarrow\;\; v = c .}
\end{equation}
A massless particle has no rest frame; the only velocity consistent with
Lorentz symmetry is the invariant speed $c$.  In Lecture~07 we will derive
electromagnetic waves from the covariant Maxwell equations and confirm that
they propagate at exactly $c$.

\subsubsection*{Coupling to the Electromagnetic Field}

To describe a particle of charge $e$ interacting with the electromagnetic
field, we add to~\eqref{eq:lec06:action-final} the only Lorentz-scalar
coupling linear in the \textbf{electromagnetic four-potential}
$A_\mu = (\phi,\,-\bvec{A})$:
\begin{equation}\label{eq:lec06:interaction}
    S_{\mathrm{int}} = -\frac{e}{c}\int A_\mu\,dx^\mu .
\end{equation}
This is a Lorentz scalar because $A_\mu$ is a covariant four-vector and
$dx^\mu$ is a contravariant four-vector; their inner product $A_\mu\,dx^\mu$
is invariant.  The total action is
\begin{equation}\label{eq:lec06:total-action}
    S = -mc\int ds - \frac{e}{c}\int A_\mu\,dx^\mu .
\end{equation}
Writing explicitly in terms of coordinate time ($\dot{x}^\mu = (c,\bvec{v})$):
\begin{equation}\label{eq:lec06:lagrangian-em}
    L = -mc^2\sqrt{1-\beta^2}
        \underbrace{-\frac{e}{c}A_\mu\,\dot{x}^\mu}_{= -e\phi + \frac{e}{c}\bvec{A}\cdot\bvec{v}}
      = -mc^2\sqrt{1-\beta^2} - e\phi + \frac{e}{c}\,\bvec{A}\cdot\bvec{v} ,
\end{equation}
which is the standard \textbf{minimal-coupling Lagrangian} for a charged
relativistic particle.

\subsubsection*{Covariant Equation of Motion: the Relativistic Lorentz Force}

Varying the total action~\eqref{eq:lec06:total-action} with respect to the
worldline $x^\mu(\lambda)$, the free-particle part yields the inertial term
$m\,du^\mu/d\tau$.  The variation of the interaction term generates the
antisymmetric combination
\begin{equation}\label{eq:lec06:field-tensor}
    F_{\mu\nu} \equiv \partial_\mu A_\nu - \partial_\nu A_\mu ,
\end{equation}
the \textbf{electromagnetic field tensor}.  ($F_{0i}$ encodes $\bvec{E}$;
$F_{ij}$ encodes $\bvec{B}$.)  Combining both contributions, the
Euler--Lagrange equations take the manifestly covariant form:
\begin{equation}\label{eq:lec06:lorentz-force}
    \boxed{m\,\frac{du^\mu}{d\tau} = \frac{e}{c}\,F^{\mu}_{\ \nu}\,u^\nu}
\end{equation}
where $u^\mu = dx^\mu/d\tau$ is the four-velocity.  This is the
\textbf{relativistic Lorentz force law}.  In the non-relativistic limit
it reduces to the familiar $m\,d\bvec{v}/dt = e\bigl(\bvec{E}
+ \bvec{v}/c\times\bvec{B}\bigr)$.

% ==============================================================================
\subsection*{Summary and Looking Ahead}
% ==============================================================================

\begin{enumerate}
    \item \textbf{Contravariant vs.\ covariant:} $V^\mu$ transforms with
      $\Lambda$; $V_\mu = \eta_{\mu\nu}V^\nu$ transforms with $\Lambda^{-1}$.
      Index position is the sole distinguishing notation.
    \item \textbf{Covariant transformation matrix:}
      $\tilde\Lambda_\mu^{\ \sigma} = \eta_{\mu\nu}\Lambda^\nu_{\ \rho}\eta^{\rho\sigma}
      = (\Lambda^{-1})^\sigma_{\ \mu}$, a direct consequence of
      $\Lambda^T\eta\Lambda = \eta$.
    \item \textbf{Mixed tensors:} index position matters;
      $T^\mu_{\ \nu} \neq T_{\ \nu}^\mu$ unless the tensor is symmetric.
    \item \textbf{Minkowski metric:} $\eta'_{\mu\nu} = \eta_{\mu\nu}$ in
      every inertial frame — a constant rank-2 covariant tensor.
    \item \textbf{Contraction:} summing one upper and one lower index reduces
      rank by 2 and yields a Lorentz-invariant object.
    \item \textbf{Modern approach:} symmetries (Lorentz + translations) fix
      the action of a free particle essentially uniquely as $S=-mc\int ds$.
      The minus sign converts the maximum of $\int ds$ to a least-action minimum.
    \item \textbf{Relativistic momentum and energy:}
      $\vec{p} = \gamma m\vec{v}$,\quad $E = \gamma mc^2$,\quad
      $E^2 = p^2c^2 + m^2c^4$.
    \item \textbf{Massless limit:} $m = 0 \Rightarrow E = pc \Rightarrow v = c$.
    \item \textbf{Electromagnetic coupling:}
      $S_{\mathrm{int}} = -(e/c)\!\int\!A_\mu\,dx^\mu$; equation of motion is
      $m\,du^\mu/d\tau = (e/c)F^{\mu}_{\ \nu}u^\nu$.
\end{enumerate}

In the next lecture we will write Maxwell's equations in manifestly covariant
form using $F^{\mu\nu}$, introduce the \textbf{energy-momentum tensor}
$T^{\mu\nu}$ of the electromagnetic field, and connect Noether's theorem to
the conservation of energy and momentum in the field.
